\chapter{نظرية الكثافة الصفرية: الإطار العام}
\personindex{ريمان، برنهارد}
\symbolindex{دالة الكثافة الصفرية \(N(\sigma,T)\)}

\section{نطاق الفصل وحالته}

\noindent\textbf{حالة الفصل:} \textenglish{AUTHORED-DRAFT / NON-CITABLE}.
هذا الفصل الأول من المجلد الخامس، ويقدّم الإطار المعياري لدالة الكثافة
\(N(\sigma,T)\) وخواصها الأولية فقط. لا يثبت هذا الفصل أي حدّ كثافة عددي
(تلك مؤجَّلة إلى الفصل التاسع والخمسين)، ولا يقترب من فرضية ريمان بأي صورة غير
مشروطة. اجتاز مسار حوكمة كامل قبل التأليف: سجل أدلة، وخريطة برهان، وثلاث
جولات مراجعة مستقلة كشفت وأصلحت خطأين رياضيين حقيقيين قبل كتابة هذا
النص. ثم اجتاز جولتَي مراجعة مستقلة بعد التأليف، كشفت أولاهما خطأين
رياضيين إضافيين في متن الفصل نفسه (اتجاه متراجحة معكوس، وتعميم غير
مشروع على الارتفاعات الصغيرة) وصُحِّحا. يبقى هذا النص بحاجة إلى
\textbf{اعتماد مالك} قبل أن يصبح قابلاً للاستشهاد؛ وحتى ذلك الحين لا
تُفعَّل أي من نتائجه ولا يُحال إليها من فصل آخر.

\subsection{المتطلبات السابقة}

يفترض الفصل:
\begin{itemize}
  \item تعريف \(N(T)\) وصيغة ريمان–فون مانغولت التقاربية لها (الفصل السادس).
  \item تناظر الأصفار غير البديهية \(\rho\leftrightarrow1-\overline\rho\)
  واحتواءها في الشريط الحرج المغلق (الفصل السادس، \textenglish{\texttt{ANT-PROP-06-01}}).
  \item المنطقة الكلاسيكية الخالية من الأصفار الخاصة بزيتا (الفصل السادس،
  \textenglish{\texttt{ANT-THM-06-09}}، مُستشهَد بها).
\end{itemize}

\subsection{نواتج التعلم}

بعد إتمام الفصل ينبغي أن يستطيع القارئ:
\begin{enumerate}
  \item تعريف \(N(\sigma,T)\) بدقة، وتمييزها عن \(N(T)\).
  \item إثبات رتابتها في \(\sigma\) والمتراجحة الأساسية التي تربطها بـ\(N(T)\).
  \item شرح لماذا مساواة هذه المتراجحة تكافئ فرضية ريمان مقصورة على ارتفاع
  معيّن، لا نتيجة أولية.
  \item استنتاج اختفاء \(N(\sigma,T)\) داخل منطقة خالية من الأصفار.
  \item صياغة حدسية الكثافة القياسية دون الخلط بينها وبين نتيجة مثبتة.
\end{enumerate}

\section{تعريف الكثافة}
\symbolindex{تعريف الكثافة}

بينما تعدّ \(N(T)\) كل الأصفار غير البديهية حتى الارتفاع \(T\) بلا تمييز،
تقيس الدالة الآتية عدد الأصفار البعيدة عن الخط الحرج جهة اليمين — وهي
الكمية الصحيحة لصياغة أي حدٍّ يقول إن أغلب الأصفار قريبة من \(\Re(s)=1/2\).

\begin{definition}[الكثافة الصفرية]
\label{def:zero-density}
\resultid{ANT-DEF-57-01}
لكل \(\sigma\in[1/2,1]\) وكل \(T\ge2\)، نضع
\[
N(\sigma,T)
=
\#\bigl\{
\rho=\beta+i\gamma
:
\zeta(\rho)=0,\
\beta\ge\sigma,\
0<\gamma\le T
\bigr\},
\]
حيث العدّ بالتكرار (رتبة كل صفر). لا نعرّف \(N(\sigma,T)\) خارج
\([1/2,1]\)؛ فهذا النطاق هو الوحيد ذو المعنى العملي لأغراض نظرية
الكثافة، مطابقةً لاصطلاح Ivić وMontgomery.
\end{definition}

\begin{remark}[نقطة الاتصال بـ\(N(T)\)]
عند \(\sigma=1/2\) لا تساوي \(N(1/2,T)\) بالضرورة \(N(T)\)؛ العلاقة
الدقيقة بينهما هي موضوع القضية التالية.
\end{remark}

\section{المتراجحة الأساسية وعلاقتها بفرضية ريمان}
\symbolindex{المتراجحة الأساسية وعلاقتها بفرضية ريمان}

\theoremindex{المتراجحة الأساسية للكثافة عند نصف}
\begin{lemma}[المتراجحة الأساسية]
\label{lem:density-half-inequality}
\resultid{ANT-LEM-57-01}
\provedhere
لكل \(T\ge2\):
\[
N(1/2,T)\le N(T).
\]
\end{lemma}

\begin{proof}
كل صفر يُحسب في \(N(1/2,T)\) يحقق \(\beta\ge1/2\)، وهو بالتحديد صفر
غير بديهي بارتفاع لا يتجاوز \(T\)، فهو محسوب أيضًا في \(N(T)\). إذن
مجموعة الأصفار المحسوبة في \(N(1/2,T)\) مجموعة جزئية من مجموعة الأصفار
المحسوبة في \(N(T)\)، والمتراجحة تنتج من مقارنة عددي مجموعتين إحداهما
جزء من الأخرى.
\end{proof}

\begin{remark}[التحذير الجوهري: المساواة تكافئ فرضية ريمان]
\label{rem:density-half-equality-is-rh}
\resultid{ANT-PRIN-57-01}
\noindent\textenglish{METHODOLOGICAL-PRINCIPLE / INFERENCE-GUARDED}.
لا تثبت هذه اللمّة، ولا يمكن أن تثبت بأدوات أولية، المساواة
\(N(1/2,T)=N(T)\). فمن تناظر الأصفار \(\rho\leftrightarrow1-\overline\rho\)
(\textenglish{\texttt{ANT-PROP-06-01}})، يأتي كل صفر بجزء حقيقي
\(\beta<1/2\) مصحوبًا بصفر "مرآتي" بجزء حقيقي \(1-\beta>1/2\) وبالارتفاع
\(\gamma\) نفسه. لذلك
\[
N(T)=N(1/2,T)+\#\{\rho:\beta<1/2,\ 0<\gamma\le T\},
\]
والمساواة \(N(1/2,T)=N(T)\) تصح إذا وفقط إذا لم يوجد أي صفر بـ
\(\beta<1/2\) حتى الارتفاع \(T\) — وبالتناظر نفسه هذا يفرض عدم وجود أي
صفر بـ\(\beta>1/2\) أيضًا، أي أن كل صفر حتى الارتفاع \(T\) يقع بالضبط
على \(\Re(s)=1/2\). هذه فرضية ريمان مقصورة على الارتفاع \(T\)، وهي مسألة
مفتوحة (الفصل السادس). لا يستعمل هذا الفصل ولا أي فصل لاحق في هذا
المجلد هذه المساواة بوصفها حقيقة مثبتة أو مفترضة.
\end{remark}

\section{الرتابة}
\symbolindex{الرتابة}

\theoremindex{رتابة الكثافة}
\begin{lemma}[رتابة الكثافة]
\label{lem:density-monotone}
\resultid{ANT-LEM-57-02}
\provedhere
الدالة \(\sigma\mapsto N(\sigma,T)\) غير متزايدة على \([1/2,1]\): لكل
\(1/2\le\sigma_1<\sigma_2\le1\)،
\[
N(\sigma_2,T)\le N(\sigma_1,T).
\]
\end{lemma}

\begin{proof}
إذا كان \(\beta\ge\sigma_2\) فإن \(\beta\ge\sigma_1\) لأن
\(\sigma_1<\sigma_2\). إذن مجموعة الأصفار المحسوبة في \(N(\sigma_2,T)\)
مجموعة جزئية من مجموعة الأصفار المحسوبة في \(N(\sigma_1,T)\)، وينتج
الحكم من مقارنة عددي المجموعتين.
\end{proof}

\section{الاختفاء داخل منطقة خالية من الأصفار}
\symbolindex{الاختفاء داخل منطقة خالية من الأصفار}

تربط القضية التالية دالة الكثافة بالمنطقة الكلاسيكية الخالية من الأصفار
الخاصة بزيتا، المُثبتة (بوصفها نتيجة مُستشهَدًا بها) في الفصل السادس.

\theoremindex{اختفاء الكثافة في منطقة خالية}
\begin{proposition}[اختفاء الكثافة في منطقة خالية، بارتفاع كافٍ]
\label{prop:density-vanishes-in-zero-free-region}
\resultid{ANT-PROP-57-01}
\provedhere
ليكن \(c_0>0\) الثابت المطلق من \textenglish{\texttt{ANT-THM-06-09}}
(المنطقة الكلاسيكية الخالية من الأصفار الخاصة بزيتا). لكل \(T\ge3\)
وكل
\[
\sigma\ge1-\frac{c_0}{\log T},
\]
لا يوجد صفر \(\rho=\beta+i\gamma\) لِـ\(\zeta\) بحيث \(\beta\ge\sigma\)
و\(3\le\gamma\le T\).
\end{proposition}

\begin{proof}
من \textenglish{\texttt{ANT-THM-06-09}}، لا يوجد صفر \(\rho=\beta+i\gamma\)
بحيث \(|\gamma|\ge3\) و\(\beta\ge1-c_0/\log|\gamma|\). الدالة
\(u\mapsto1-c_0/\log u\) متزايدة على \(u\ge3\) (لأن \(\log u\) متزايدة
وموجبة هناك، فـ\(c_0/\log u\) متناقصة). إذن كل \(\gamma\) يحقق
\(3\le\gamma\le T\) يعطي \(\log\gamma\le\log T\)، ومن ثم
\[
1-\frac{c_0}{\log\gamma}\le1-\frac{c_0}{\log T}.
\]
فأي صفر بارتفاع \(3\le\gamma\le T\) يحقق
\[
\beta<1-\frac{c_0}{\log\gamma}\le1-\frac{c_0}{\log T}\le\sigma,
\]
أي \(\beta<\sigma\). إذن لا يوجد صفر من هذا النوع يحقق \(\beta\ge\sigma\).
\end{proof}

\begin{remark}[الفجوة عند الارتفاعات الصغيرة، مؤجَّلة صراحةً]
\label{rem:density-small-height-gap}
القضية أعلاه، كما بُرهنت، تستبعد فقط الأصفار بارتفاع \(\gamma\ge3\).
أما الأصفار المحتملة بارتفاع \(0<\gamma<3\) فلا تغطيها
\textenglish{\texttt{ANT-THM-06-09}} (المقصورة صراحةً على \(|t|\ge3\))؛
استبعادها عند \(\sigma\) قريبة من \(1\) يتطلب مدخلًا إضافيًا — إما
استشهادًا صريحًا بغياب مثل هذه الأصفار (متوفر عدديًّا في الأدبيات
القياسية، إذ يقع أول صفر غير بديهي معروف عند ارتفاع أكبر من \(14\)،
لكن هذه الحقيقة لم تُدرَج بعد في الموسوعة بمعرّف قابل للاستشهاد) أو
حجة تحليلية مستقلة لم تُقدَّم هنا. لذلك \(N(\sigma,T)=0\) بالمعنى الكامل
(لكل \(0<\gamma\le T\)) \textbf{ليست} نتيجة هذا الفصل؛ ما أُثبت هو غياب هذه
الأصفار بين تلك ذات الارتفاع \(\gamma\ge3\) فقط. انظر «ما الذي لم يثبته
هذا الفصل؟».
\end{remark}

\begin{remark}[ملاحظة عدم-الاستنتاج، مكرَّرة من الفصل السادس]
مطابقةً لتحذير الفصل السادس صراحةً: لا تُستنتج
\textenglish{\texttt{ANT-THM-06-09}} من مبرهنة الفصل الحادي عشر
(\textenglish{\texttt{ANT-THM-11-01}}) بوضع الموصل \(q=1\)؛ تلك الأخيرة
مصاغة حصرًا لشخصيات ديريشليه البدائية غير الرئيسية بموصل \(q\ge3\)، ولا
تشمل حالة زيتا. القضية أعلاه تعتمد \textenglish{\texttt{ANT-THM-06-09}}
وحدها.
\end{remark}

\section{نتيجة: حدّ أعلى موحَّد}
\symbolindex{نتيجة: حدّ أعلى موحَّد}

\theoremindex{الحدّ الأعلى الموحَّد للكثافة}
\begin{corollary}[الحدّ الأعلى الموحَّد]
\label{cor:density-upper-bound-uniform}
\resultid{ANT-COR-57-01}
\provedhere
لكل \(\sigma\in[1/2,1]\) وكل \(T\ge2\):
\[
N(\sigma,T)\le N(T).
\]
\end{corollary}

\begin{proof}
من رتابة \(N(\sigma,T)\) في \(\sigma\)
(\ref{lem:density-monotone})، \(N(\sigma,T)\le N(1/2,T)\) لكل
\(\sigma\ge1/2\). ومن المتراجحة الأساسية (\ref{lem:density-half-inequality})،
\(N(1/2,T)\le N(T)\). يعطي تسلسل المتراجحتين الحكم مباشرة.
\end{proof}

\section{حدسية الكثافة}
\symbolindex{حدسية الكثافة}

القضايا السابقة كلها متراجحات بسيطة ونتيجة اختفاء موضعي. أما الهدف
الحقيقي لنظرية الكثافة، والذي يبرر استقلالها بفصول لاحقة، فهو تقدير
عددي صريح لـ\(N(\sigma,T)\) كدالة في \(\sigma\) و\(T\) معًا — لا مجرد
حكم بالاختفاء أو عدمه.

\begin{conjecture}[حدسية الكثافة]
\label{conj:density-hypothesis}
\resultid{ANT-CONJ-57-01}
لكل \(\sigma\in[1/2,1]\) ولكل \(\varepsilon>0\):
\[
N(\sigma,T)\ll_\varepsilon T^{2(1-\sigma)+\varepsilon}.
\]
\end{conjecture}

\begin{remark}
عند \(\sigma=1/2\) يعطي الأس \(2(1-\sigma)=1\)، متوافقًا مع الحد التافه
\(N(1/2,T)\le N(T)\ll T\log T\) من صيغة عدّ الأصفار. عند \(\sigma=1\)
يعطي الأس \(0\)، متوافقًا مع اختفاء الكثافة قرب الحافة اليمنى في مناطق
خالية كافية العرض. لا يُثبت هذا الفصل أي حدٍّ أقوى من الحد التافه عند
\(\sigma=1/2\)؛ الحدود غير التافهة (Ingham، Huxley، وما بعدهما) موضوع
الفصل التاسع والخمسين حصرًا.
\end{remark}

\section{ما الذي لم يثبته هذا الفصل؟}
\symbolindex{ما الذي لم يثبته هذا الفصل؟}

\begin{itemize}
  \item لم يثبت أي حدّ كثافة غير تافه (\(N(\sigma,T)\ll T^{c}\) بأس
  \(c<1\) عند \(\sigma>1/2\) ثابت). هذا موضوع الفصل التاسع والخمسين.
  \item لم يثبت أي شكل من فرضية ريمان، ولا اقترب من إثباتها؛ المتراجحة
  الأساسية (\ref{lem:density-half-inequality}) متراجحة أولية لا علاقة
  مباشرة لها بإثبات أو دحض الفرضية.
  \item لم يستعمل الفصل حدسية الكثافة (\ref{conj:density-hypothesis})
  في إثبات أي نتيجة؛ ذُكرت بوصفها هدفًا للفصول اللاحقة فقط.
  \item لم تُستبعد أصفار محتملة بارتفاع صغير \(0<\gamma<3\) من قضية
  الاختفاء (\ref{prop:density-vanishes-in-zero-free-region})؛ القضية
  مبرهنة فقط للأصفار بارتفاع \(\gamma\ge3\)، لأن
  \textenglish{\texttt{ANT-THM-06-09}} مقصورة على \(|t|\ge3\) (انظر
  الملاحظة~\ref{rem:density-small-height-gap}). إغلاق هذه الفجوة
  بالكامل — تحليليًّا أو باستشهاد عددي قابل للاستشهاد — مؤجَّل.
\end{itemize}

\section{تمارين}
\symbolindex{تمارين}

\begin{enumerate}
  \item بيّن أن \(N(\sigma,T)\) عدد صحيح غير سالب لكل \(\sigma,T\) في
  نطاق التعريف.
  \item بيّن أولًا أن حجة تضمّن المجموعات في برهان المتراجحة الأساسية
  (\ref{lem:density-half-inequality}) تعطي \(N(3/4,T)\le N(T)\)
  \textbf{مباشرةً}، بلا حاجة إلى الرتابة
  (\ref{lem:density-monotone}) إطلاقًا. ثم اشرح لماذا لا يعطي هذا
  البرهان — ولا الرتابة معه — أي شيء \textbf{أقوى} من هذا الحدّ التافه:
  أي معلومة إضافية يلزم إدخالها لنيل حدّ من الشكل
  \(N(3/4,T)\ll T^{c}\) بأسّ \(c<1\)؟
  \item بافتراض صحة حدسية الكثافة، استنتج (شكليًّا، دون برهان صارم هنا)
  أن نسبة الأصفار بـ\(\beta\ge3/4\) إلى إجمالي الأصفار حتى الارتفاع
  \(T\) تؤول إلى الصفر عندما \(T\to\infty\). أين بالضبط يدخل الأس
  \(2(1-\sigma)<1\) في هذا الاستنتاج؟
  \item وضّح الفرق بين "منطقة خالية من الأصفار" (نتيجة قطعية: صفر بالضبط)
  و"حدسية الكثافة" (نتيجة إحصائية: قليل جدًّا، لا صفر بالضرورة).
\end{enumerate}

\section{خلاصة الفصل}
\symbolindex{خلاصة الفصل}

قدّم هذا الفصل \(N(\sigma,T)\) بوصفها الأداة الصحيحة لقياس تركّز الأصفار
قرب الخط الحرج، وأثبت عنها ثلاث حقائق أولية فقط: متراجحة أساسية (لا
مساواة — المساواة تكافئ فرضية ريمان)، رتابة في \(\sigma\)، واختفاء داخل
منطقة خالية من الأصفار (لأصفار بارتفاع \(\gamma\ge3\) فقط) مستوردة من
الفصل السادس. صيغت حدسية الكثافة
القياسية بوصفها هدفًا مفتوحًا صراحةً، لا نتيجة. هذا الإطار الضيق هو ما
يحتاجه الفصل التاسع والخمسون ليبدأ بحدود Ingham وHuxley العددية دون إعادة بناء
أي من هذه الأساسيات.
